% Part: turing-machines
% Chapter: machines-computations
% Section: halting-states

\documentclass[../../../include/open-logic-section]{subfiles}

\begin{document}

\olfileid{tur}{mac}{hal}
\olsection{رُکّن والیاں حالتاں}

\begin{explain}
بھاویں اسی اپنی مشیناں
لئی ایہ تعریف رکھی اے
کہ اوہ صرف اُدوں رُکدیاں
نیں جدوں چلاؤن لئی
کوئی حکم نہ ہووے، ٹیورنگ
مشیناں دی عام نمائندگی
وچ اک خاص \emph{رُکّن
والی حالت}~$h$ ہوندی
اے، تے $h \in Q$ اے۔

رُکّن والی حالت دا
خیال سادا اے: جدوں
مشین اپنا کم مکمل
کر لیندی اے (انپٹ
قبول کرن لئی تیار
ہو جاوے، یا آؤٹ پٹ
لکھ چکی ہووے)، اوہ
حالت~$h$ وچ چلی جاندی
اے تے اوس حالت وچ رُک جاندی اے۔
کجھ مشیناں دیاں
رُکّن والیاں دو حالتاں
ہوندیاں نیں: اک انپٹ
قبول کرن لئی تے
دوجی انپٹ رد کرن لئی۔
\end{explain}

\begin{ex}\emph{رُکّن والیاں حالتاں}۔
ایس خیال نوں کھولن لئی
اسی جفت مشین نوں تھوڑھا
بدل کے شروع کردے آں۔
جے انپٹ دیاں لکیراں دی گنتی جفت ہووے تاں
مشین نوں حالت~$q_0$
وچ رُکّان دی بجائے
اسی اک حکم ودھا سکدے
آں جیڑا مشین نوں رُکّن
والی حالت وچ لے جاوے۔
\[
\begin{tikzpicture}[->,>=stealth',shorten >=1pt,auto,node distance=2.8cm,
                    semithick]
  \tikzstyle{every state}=[fill=none,draw=black,text=black]

  \node[initial, state]         (A)                     {$q_0$};
  \node[state]         (B) [right of=A] {$q_1$};
  \node[state]         (C) [below of=A] {$h$};

  \path (A) edge [bend left] node {\TMtrans{\TMstroke}{\TMstroke}{R}} (B)
  	    edge node {\TMtrans{\TMblank}{\TMblank}{N}} (C)
        (B) edge [loop above] node {\TMtrans{\TMblank}{\TMblank}{R}} (B)
            edge [bend left] node {\TMtrans{\TMstroke}{\TMstroke}{R}} (A);
\end{tikzpicture}
\]

ہن مثال نوں ہور ودھائیے۔
جدوں مشین نوں پتا لگدا اے
کہ انپٹ دیاں لکیراں دی گنتی طاق اے،
اوہ کدی نہیں رُکدی۔ لوپ والے حکم نوں
رد کرن والی حالت~$r$
تک جان والے حکم نال
بدل کے اسی مشین وچ
اک \emph{رد} والی حالت
شامل کر سکدے آں۔
\[
\begin{tikzpicture}[->,>=stealth',shorten >=1pt,auto,node distance=2.8cm,
                    semithick]
  \tikzstyle{every state}=[fill=none,draw=black,text=black]

  \node[initial,state]         (A)                     {$q_0$};
  \node[state]         (B) [right of=A] {$q_1$};
  \node[state]         (C) [below of=A] {$h$};
  \node[state]         (D) [below of=B] {$r$};

  \path (A) edge [bend left] node {\TMtrans{\TMstroke}{\TMstroke}{R}} (B)
  	    edge node {\TMtrans{\TMblank}{\TMblank}{N}} (C)
        (B) edge node {\TMtrans{\TMblank}{\TMblank}{N}} (D)
            edge [bend left] node {\TMtrans{\TMstroke}{\TMstroke}{R}} (A);
\end{tikzpicture}
\]
\end{ex}

\begin{explain}
ایہو جہیاں صورتاں وچ
خاص رُکّن والی حالت
شامل کرنی فائدہ مند
ہو سکدی اے، کیوں جے
ایس نال صاف پتا لگدا
اے کہ مشین کدوں کسے
انپٹ نوں قبول یا رد
کردی اے۔ پر ایہ گل
یاد رکھنی ضروری اے
کہ خاص رُکّن والی
حالت ودھاؤن نال مشین
دی حساب کرن دی طاقت
نہیں ودھدی۔ ایس طرح
رُکّن دے گھٹ رسمی
خیال دے وی اپنے فائدے
نیں۔ ایس باب وچ ہن
تک ورتی گئی رُکّن
دی تعریف نال
\emph{رُکّن دے مسئلے}
بارے ثبوت سمجھنا تے
وکھانا سوکھا ہو جاندا اے۔ ایس
لئی اسی اپنی اصل
تعریف ہی ورتدے رہندے
آں۔
\end{explain}

\end{document}
